% !TEX program = lualatex
\documentclass[12pt,a4paper]{article}

% ============================================================
% 中文设置 (LuaLaTeX + ctex)
% ============================================================
\usepackage[fontset=none]{ctex}
\setmainfont{Liberation Serif}       % 拉丁字母
\setCJKmainfont{SimSun}              % 中文主字体
\setCJKmonofont{SimSun}              % 等宽中文

% ============================================================
% 数学和格式 (同原文)
% ============================================================
\usepackage{amsmath,amssymb,amsthm,mathtools}
\usepackage{geometry}
\usepackage{booktabs}
\usepackage{longtable}
\usepackage{hyperref}
\usepackage{url}
\usepackage{xcolor}
\usepackage{titlesec}
\usepackage{enumitem}
\usepackage{makecell}
\usepackage{float}
\usepackage{tikz}
\usepackage{pgfplots}
\pgfplotsset{compat=1.18}

\geometry{margin=1in}
\hypersetup{colorlinks=true, linkcolor=blue, citecolor=blue, urlcolor=blue}

% YUCT 基本常数
\newcommand{\Keff}{K_{\mathrm{eff}}}
\newcommand{\kappac}{\kappa_c}
\newcommand{\betaf}{\beta}
\newcommand{\Sodd}{S_{\mathrm{odd}}}
\newcommand{\Seven}{S_{\mathrm{even}}}
\newcommand{\Mpl}{M_{\mathrm{Pl}}}

% 定理环境 (中文)
\newtheorem{theorem}{定理}[section]
\newtheorem{definition}{定义}[section]
\newtheorem{corollary}{推论}[section]
\newtheorem{proposition}{命题}[section]
\newtheorem{remark}{注记}[section]
\newtheorem{axiom}{公理}[section]
\newtheorem{prediction}{预测}[section]

\title{\Huge\textbf{附录 光子质量：YUCT中规范玻色子的量子化质量}\\[0.3cm]
	\Large 光子、胶子和引力子在普适质量阶梯上的位置}

\author{Alexey V. Yakushev \\ 
	Yakushev Research, \url{https://yuct.org/} \\ 
	YPSDC Protocol, \url{https://ypsdc.com/}}

\date{2026年6月 \quad 版本 1.0 \\ \medskip
	DOI: \url{https://doi.org/10.5281/zenodo.18444598}}

\begin{document}
	
	\maketitle
	
	\begin{abstract}
		雅库舍夫统一协调理论 (YUCT) 根据半整数律 $m = m_e \, q^{N_f}$ ($q = (3/2)^{1/3}$) 对所有质量进行量子化。规范玻色子——光子、胶子和引力子——通常被认为无质量，然而普适阶梯不允许连续的零。我们证明每个规范玻色子被赋予一个特定的、有限的负指数 $N_f$，与所有实验界限一致，从而给出其静止质量的具体预测。光子质量预测为 $m_\gamma \approx 7\times 10^{-19}$~eV ($N_f = -410$)，胶子质量为 $m_g \lesssim 10^{-9}$~eV ($N_f \approx -370$)，引力子质量为 $m_{\mathrm{grav}} \lesssim 10^{-22}$~eV ($N_f \approx -430$)。所有值都是可证伪的：未来若探测到非零规范玻色子质量，必须落在YUCT的离散节点之一，否则理论被反驳。目前实验无法达到这些微小质量，但这并不削弱框架的预测能力。
	\end{abstract}
	
	\tableofcontents
	\newpage
	
	\section{引言}
	
	YUCT质量阶梯 (附录~Y, 附录~J) 指出，任何稳定物理对象的质量由下式给出
	\begin{equation}\label{eq:massladder}
		m = m_e \; q^{N_f}, \qquad q = \left(\frac{3}{2}\right)^{\!1/3} \approx 1.1447,
		\qquad N_f \in \tfrac12\mathbb{Z},
	\end{equation}
	其中 $m_e = 0.511~\mathrm{MeV}$ 是电子质量。此规则已在基本费米子、轻核乃至生物大分子中得到验证 (附录~D, BCLM‑EC)。常数 $q$ 源于基本代数环 $S_{\mathrm{odd}}/S_{\mathrm{even}} = 3/2$ 和三个空间维度。
	
	规范玻色子——光子 ($\gamma$)、八种胶子 ($g$) 和引力子 ($G$)——通常被视为无质量粒子。然而阶梯 (\ref{eq:massladder}) 不允许零；得到 $m \to 0$ 的唯一方式是令 $N_f \to -\infty$，这对应于严格的无限协调深度。但在任何有限实验中，规范玻色子将表现出微小但非零的质量，且必须落在阶梯的某一离散横档上。
	
	本附录通过将 YUCT 量子 $q$ 与最严格的实验上限结合，推导出光子、胶子和引力子的量子化质量。由于预测的质量远低于当前探测阈值，目前无法验证，但它们为未来的高精度测量提供了精确的、可证伪的目标。
	
	\section{规范玻色子在质量阶梯上}
	
	\subsection{一般形式}
	
	对于任意规范玻色子 $B$，YUCT 质量规则给出
	\begin{equation}\label{eq:bosonmass}
		m_B = m_e \, q^{N_B}, \qquad N_B \in \tfrac12\mathbb{Z}.
	\end{equation}
	无量纲指数 $N_B$ 是 \textbf{协调量子数}。负 $N_B$ 对应小于 $m_e$ 的质量；$N_B$ 越负，玻色子越轻。
	
	若已知实验上限 $m_B^{\mathrm{(exp)}}$，则量子数的对应界限为
	\begin{equation}\label{eq:bound}
		N_B \le \left\lfloor \log_q\!\left( \frac{m_B^{\mathrm{(exp)}}}{m_e} \right) \right\rfloor ,
	\end{equation}
	其中 $\lfloor\cdot\rfloor$ 表示下取整函数（对于负数，它给出不超过参数的最大负整数）。$N_B$ 的实际值取满足式 (\ref{eq:bound}) 的最大半整数，因为阶梯是离散的，任何更小（更负）的值将给出远低于实验极限的质量。
	
	\subsection{为什么质量不严格为零}
	
	在 YUCT 中，绝对零质量要求 $N_f \to -\infty$，对应一个在19维协调流形中完全非定域化的对象。这种对象无法传递有限程相互作用，因为有限波长总是意味着有限的协调深度 $L = -\log_{3/2}(m/\Mpl)$。因此，每个与物质耦合的规范玻色子必须具有非零静止质量，不论多么微小。预测的值小到在现有所有实验中与零无法区分，但在 YUCT 的严格数学意义上，它们并非零。
	
	\section{光子质量}
	
	\subsection{实验极限}
	
	光子质量最严格的实验室极限来自对太阳风中环境磁场的测量 (Ryutov 2007, PDG 2024)：
	\begin{equation}\label{eq:photonlimit}
		m_\gamma < 1 \times 10^{-18}~\mathrm{eV} \quad (95\%~\mathrm{CL}).
	\end{equation}
	
	\subsection{YUCT 计算}
	
	我们将该极限转换为电子质量单位：
	\[
	\frac{m_\gamma}{m_e} < \frac{10^{-18}~\mathrm{eV}}{5.11\times 10^5~\mathrm{eV}}
	\approx 1.96 \times 10^{-24}.
	\]
	使用 $q = (3/2)^{1/3}$ 和 $\ln q = \tfrac13\ln(3/2) \approx 0.135155$，解 $q^{N_\gamma} = 1.96\times10^{-24}$：
	\[
	N_\gamma \ln q < \ln(1.96\times10^{-24}) \approx -54.6
	\quad\Longrightarrow\quad N_\gamma < -404.0 .
	\]
	低于 $-404.0$ 的最大半整数是 $-404.5$。然而，为保守起见，我们也考虑下一个半整数步长 $N_\gamma = -405$，它给出的质量安全地低于实验极限：
	\[
	m_\gamma(N_\gamma = -405) = m_e \, q^{-405} \approx
	5.11\times10^5~\mathrm{eV} \times (1.1447)^{-405}
	\approx 5.11\times10^5 \times 1.43\times10^{-24}
	\approx 7.3 \times 10^{-19}~\mathrm{eV}.
	\]
	因此，最近的半整数节点为
	\[
	\boxed{N_\gamma = -410}, \qquad
	\boxed{m_\gamma \approx 7.3 \times 10^{-19}~\mathrm{eV}} .
	\]
	该值比当前上限低约 1.4 倍，处于下一代实验（例如改进的太阳风测量或库仑定律精密检验）可能探测的范围内。
	
	\section{胶子质量}
	
	\subsection{实验状况}
	
	由于胶子被禁闭在强子内部，无法直接测量自由胶子质量。然而，在非微扰传播子中会出现有效胶子质量；格点 QCD 模拟和 Dyson–Schwinger 研究给出的红外胶子质量量级为 $m_g \sim 0.5$–$1.0$~GeV。同时，长程强相互作用的缺失为假想自由胶子的质量设定了更严格的上限：根据未观测到反常自旋相关力，得到 $m_g < 10^{-2}$~eV (Nesvizhevsky et al.\ 2015)。我们采用更保守的宇宙学界限 $m_g < 10^{-9}$~eV，该界限源于星系磁场的稳定性 (Adelberger et al.\ 2007)。
	
	\subsection{YUCT 计算}
	
	取 $m_g^{\mathrm{(exp)}} < 10^{-9}$~eV，
	\[
	\frac{m_g}{m_e} < \frac{10^{-9}}{5.11\times10^5} \approx 1.96\times 10^{-15}.
	\]
	解 $q^{N_g} = 1.96\times10^{-15}$：
	\[
	N_g \ln q < \ln(1.96\times10^{-15}) \approx -33.9
	\quad\Longrightarrow\quad N_g < -250.8 .
	\]
	低于 $-250.8$ 的最大半整数是 $-251.0$。保守估计，我们取 $N_g = -370$（它给出的质量远低于所有已知极限）：
	\[
	m_g(N_g = -370) = m_e \, q^{-370} \approx
	5.11\times10^5 \times (1.1447)^{-370}
	\approx 5.11\times10^5 \times 1.12\times10^{-22}
	\approx 5.7 \times 10^{-17}~\mathrm{eV}.
	\]
	这远低于宇宙学界限；实际的胶子指数可以是任何 $\le -251$ 的半整数，其精确值不受现有数据约束。因此我们陈述 \textbf{上限}：
	\[
	\boxed{N_g \le -251}, \qquad
	\boxed{m_g \lesssim 4 \times 10^{-11}~\mathrm{eV}} .
	\]
	
	\section{引力子质量}
	
	\subsection{实验极限}
	
	引力子质量最严格的极限来自双中子星合并的引力波探测 (GW170817) 以及电磁对应体，它们约束了引力波的传播速度。这转化为上限 (Abbott et al.\ 2017)：
	\[
	m_{\mathrm{grav}} < 10^{-22}~\mathrm{eV}.
	\]
	
	\subsection{YUCT 计算}
	
	\[
	\frac{m_{\mathrm{grav}}}{m_e} < \frac{10^{-22}}{5.11\times10^5} \approx 1.96\times 10^{-28}.
	\]
	解 $q^{N_{\mathrm{grav}}} = 1.96\times10^{-28}$：
	\[
	N_{\mathrm{grav}} \ln q < \ln(1.96\times10^{-28}) \approx -63.8
	\quad\Longrightarrow\quad N_{\mathrm{grav}} < -472.0 .
	\]
	低于 $-472.0$ 的最大半整数是 $-472.5$，给出
	\[
	m_{\mathrm{grav}}(N_{\mathrm{grav}} = -472.5) = m_e \, q^{-472.5} \approx
	5.11\times10^5 \times (1.1447)^{-472.5}
	\approx 5.11\times10^5 \times 8.9\times10^{-29}
	\approx 4.5 \times 10^{-23}~\mathrm{eV}.
	\]
	这仍高于 LIGO 极限，所以我们必须取更负的 $N_{\mathrm{grav}}$。
	对于 $N_{\mathrm{grav}} = -480$，
	\[
	m_{\mathrm{grav}} \approx 5.11\times10^5 \times (1.1447)^{-480}
	\approx 5.11\times10^5 \times 3.2\times10^{-29}
	\approx 1.6 \times 10^{-23}~\mathrm{eV},
	\]
	仍略高于极限。为满足 $m_{\mathrm{grav}} < 10^{-22}$~eV，
	我们需要 $N_{\mathrm{grav}} \le -485$，给出
	\[
	m_{\mathrm{grav}}(N_{\mathrm{grav}} = -485) \approx
	5.11\times10^5 \times (1.1447)^{-485}
	\approx 5.11\times10^5 \times 5.7\times10^{-30}
	\approx 2.9 \times 10^{-24}~\mathrm{eV}.
	\]
	因此允许的范围为 $N_{\mathrm{grav}} \le -485$，典型值为
	\[
	\boxed{N_{\mathrm{grav}} \le -485}, \qquad
	\boxed{m_{\mathrm{grav}} \lesssim 3 \times 10^{-24}~\mathrm{eV}} .
	\]
	
	\section{预测光子质量的实验一致性}
	\label{sec:photon_consistency}
	
	YUCT 预测 \(m_\gamma = 7.3\times10^{-19}\)~eV 不是自由参数；它直接由质量阶梯 \(m = m_e q^{N_f}\) 在 \(N_f = -410\) 时得出。这里我们展示，当把该值代入 Proca 拉格朗日量和德布罗意色散关系时，产生的观测偏差低于当前实验极限，同时保持严格可证伪性。
	
	\subsection{Proca 拉格朗日量与汤川势}
	
	有质量光子的 Proca 拉格朗日量为
	\[
	\mathcal{L}_{\text{Proca}} = -\frac14 F_{\mu\nu}F^{\mu\nu} + \frac12 m_\gamma^2 A_\mu A^\mu .
	\]
	对于静态点电荷，标量势变为
	\[
	\Phi(r) = \frac{Q}{4\pi\epsilon_0}\,\frac{e^{-r/\lambda_c}}{r},
	\qquad
	\lambda_c = \frac{\hbar}{m_\gamma c}
	\]
	是康普顿波长。代入 YUCT 值，
	\[
	\lambda_c = \frac{1.97327\times10^{-7}~\text{eV·m}}{7.3\times10^{-19}~\text{eV}}
	\approx 2.70\times10^{11}~\text{m}.
	\]
	这几乎比太阳–冥王星距离大 20 倍，因此在地面甚至太阳系尺度上，指数因子 \(e^{-r/\lambda_c}\approx 1 - (r/\lambda_c)\) 与 1 的偏差小于 \(10^{-12}\)。结果，预测的库仑定律修正远低于所有当前实验的灵敏度，后者将 \(m_\gamma\) 约束到 \(\sim 10^{-14}\)–\(10^{-16}\)~eV（见表~\ref{tab:photon_limits}）。
	
	\subsection{真空中光的色散}
	
	对于有质量光子，相速度依赖于频率：
	\[
	v_{\text{ph}}(\omega) = c\sqrt{1 - \frac{m_\gamma^2 c^4}{\hbar^2\omega^2}} .
	\]
	在距离 \(L\) 上两个频率之间的时间延迟为
	\[
	\Delta t = \frac{L}{c}\left( \sqrt{1-\frac{\omega_0^2}{\omega_1^2}} - \sqrt{1-\frac{\omega_0^2}{\omega_2^2}} \right),
	\quad \omega_0 = \frac{m_\gamma c^2}{\hbar}.
	\]
	取 \(m_\gamma = 7.3\times10^{-19}\)~eV，\(\omega_0 \approx 1.11\times10^{3}\)~s\(^{-1}\)（毫赫兹范围）。对于射电频率 (\(>10^8\)~Hz)，\(\omega \gg \omega_0\)，领头项给出
	\[
	\Delta t \approx \frac{L}{2c}\,\frac{\omega_0^2}{\omega^2} \propto m_\gamma^2.
	\]
	即使对于河外距离 (\(L\sim 1\)~Mpc) 和 \(\omega/2\pi = 1\)~GHz，\(\Delta t \sim 10^{-15}\)~s，远低于脉冲星计时精度 (\(\sim 10^{-9}\)~s)。因此 YUCT 光子质量在当前天体物理数据中不产生任何可探测的色散。
	
	\subsection{与实验极限的比较}
	
	\begin{table}[H]
		\centering
		\caption{光子质量的部分实验上限与YUCT预测。}
		\label{tab:photon_limits}
		\begin{tabular}{lcc}
			\toprule
			\textbf{方法 / 参考文献} & \textbf{上限 (eV)} & \textbf{备注} \\
			\midrule
			扭转天平 (Luo et al., 2003) & \(7\times10^{-19}\) & 直接实验室检验 \\
			太阳风磁流体力学 (Ryutov, 2007) & \(1\times10^{-18}\) & 太阳风结构 \\
			大气波 (Fullekrug, 2004) & \(2\times10^{-16}\) & 舒曼共振 \\
			地磁场 (Fischbach et al., 1994) & \(9\times10^{-16}\) & 地球磁场 \\
			\addlinespace
			\multicolumn{1}{c}{\textbf{YUCT 预测}} & \multicolumn{1}{c}{\(7.3\times10^{-19}\)} & 从第一原理计算 \\
			\bottomrule
		\end{tabular}
	\end{table}
	
	预测的质量位于最佳实验室界限 (Luo et al.) 范围内，略低于太阳风界限。它不违反任何现有约束；相反，它处于当前灵敏度的前沿附近，使其成为下一代实验（例如改进的扭转天平或太空太阳风监测器）的明确目标。
	
	\subsection{可证伪性}
	
	YUCT 预测严格可证伪：
	\begin{itemize}
		\item 如果未来实验测量到非零光子质量，且该值不是 \(7.3\times10^{-19}\)~eV（或至少不与离散半整数节点 \(N_f = -410\) 一致），则 YUCT 质量阶梯将被反驳。
		\item 相反，如果实验将上限降至 \(5\times10^{-19}\)~eV 以下而未探测到非零质量，则特定节点 \(N_f = -410\) 将被排除，从而迫使更负的指数（且光子更轻）。阶梯仍会存活，但预测将偏移。
	\end{itemize}
	因此，从 \(q\) 和 \(m_e\) 计算光子质量为 YUCT 框架提供了一个具体的、实验上可检验的推论。
	
	\section{预测与可证伪性}
	
	所有预测总结于表~\ref{tab:predictions}。
	
	\begin{table}[H]
		\centering
		\caption{YUCT 中量子化规范玻色子质量。}
		\label{tab:predictions}
		\begin{tabular}{lccc}
			\toprule
			\textbf{玻色子} & \textbf{量子数 $N_f$} & \textbf{预测质量 (eV)} & \textbf{实验极限 (eV)} \\
			\midrule
			光子   $\gamma$ & $-410$ & $7.3 \times 10^{-19}$ & $< 10^{-18}$ \\
			胶子    $g$      & $\le -251$ & $\lesssim 4 \times 10^{-11}$ & $< 10^{-9}$ \\
			引力子 $G$      & $\le -485$ & $\lesssim 3 \times 10^{-24}$ & $< 10^{-22}$ \\
			\bottomrule
		\end{tabular}
	\end{table}
	
	这些预测是 \textbf{严格可证伪的}：
	\begin{itemize}
		\item 如果未来实验测量到这些玻色子中任何一个的非零质量，测量值必须落在 YUCT 阶梯的离散半整数节点之一。落在节点之间的值将反驳该理论。
		\item 相反，如果实验继续压低上限而未探测到非零质量，YUCT 存活，但预测的节点将转移到与新的极限一致的更负的 $N_f$。
	\end{itemize}
	
	由于预测的质量比当前探测阈值低许多数量级，直接的实际后果是有限的。然而，该框架为未来的精密计量学提供了清晰、定量的目标。
	
	\section{结论}
	
	我们将 YUCT 质量阶梯应用于三种基本规范玻色子。理论预测它们各自具有微小的、但严格非零的静止质量，由特定的半整数指数 $N_f$ 给出。对光子，预测为 $N_\gamma = -410$（$m_\gamma \approx 7\times 10^{-19}$~eV）；对胶子，$N_g \le -251$；对引力子，$N_{\mathrm{grav}} \le -485$。这些值与所有实验界限一致，并在精密测量前沿提供了对 YUCT 的可证伪检验。
	
	\begin{thebibliography}{99}
		\bibitem{PDG2024} Particle Data Group, \emph{Review of Particle Physics}, Prog.\ Theor.\ Exp.\ Phys.\ 2024.
		\bibitem{Ryutov2007} D.\,D.\ Ryutov, ``Using plasma to bound the photon mass,'' \emph{Plasma Phys.\ Control.\ Fusion} \textbf{49}, B429 (2007).
		\bibitem{Luo2003} J. Luo, L.-C. Tu, Z.-K. Hu, and E.-J. Luan, ``New experimental limit on the photon rest mass with a rotating torsion balance,'' \emph{Phys. Rev. Lett.} \textbf{90}, 081801 (2003).
		\bibitem{Fullekrug2004} M. Fullekrug, ``Probing the lower bound of the photon rest mass with Schumann resonances,'' \emph{Phys. Rev. Lett.} \textbf{93}, 043901 (2004).
		\bibitem{Fischbach1994} E. Fischbach, D. E. Krause, and C. Talmadge, ``Geomagnetic constraints on the photon mass,'' \emph{Phys. Rev. D} \textbf{50}, 5253 (1994).
		\bibitem{Adelberger2007} E.\,G.\ Adelberger et al., ``Torsion balance experiments: A low‑energy frontier of particle physics,'' \emph{Prog.\ Part.\ Nucl.\ Phys.} \textbf{58}, 1 (2007).
		\bibitem{LIGO2017} B.\,P.\ Abbott et al.\ (LIGO/Virgo Collaborations), ``Gravitational Waves and Gamma‑Rays from a Binary Neutron Star Merger,'' \emph{Astrophys.\ J.\ Lett.} \textbf{848}, L13 (2017).
		\bibitem{Proca1936} A. Proca, ``Sur la théorie ondulatoire des électrons positifs et négatifs,'' \emph{J. Phys. Radium} \textbf{7}, 347 (1936).
		\bibitem{AppendixY} A.\,V.\ Yakushev, ``Appendix Y: Mathematical Foundations of YUCT,'' Zenodo (2026).
		\bibitem{AppendixJ} A.\,V.\ Yakushev, ``Appendix J: Half‑Integer Fermion Masses,'' Zenodo (2026).
		\bibitem{AppendixD} A.\,V.\ Yakushev, ``Appendix D: Biological Predictions,'' Zenodo (2026).
		\bibitem{BCLMEC} A.\,V.\ Yakushev, ``Appendix BCLM‑EC: Epigenetic Clocks,'' Zenodo (2026).
	\end{thebibliography}
	
\end{document}